%%%%%%%%%%%%%%%%%%%%%%%%%%%%%%%%%%%%%%%%%%%%%%%%%%%%%%%%%%%%%%%%%%%%%%%%%%%%%%%%
% IEEE Transactions on Power Systems Paper
% Differentiable Robust Optimization for Learning Correlated Uncertainty Sets
% in Operating Reserve Markets
%%%%%%%%%%%%%%%%%%%%%%%%%%%%%%%%%%%%%%%%%%%%%%%%%%%%%%%%%%%%%%%%%%%%%%%%%%%%%%%%

\documentclass[lettersize,journal]{IEEEtran}

%===============================================================================
% Packages
%===============================================================================
\usepackage{amsmath,amssymb,amsthm}
\usepackage{mathtools}
\usepackage{bm}
\usepackage{algorithm}
\usepackage{algorithmic}
\usepackage{booktabs}
\usepackage{multirow}
\usepackage{array}
\usepackage{cite}
\usepackage{url}
\usepackage{graphicx}
\usepackage{textcomp}

%===============================================================================
% Custom Commands
%===============================================================================
\newcommand{\R}{\mathbb{R}}
\newcommand{\E}{\mathbb{E}}
\newcommand{\Prob}{\mathbb{P}}
\newcommand{\cU}{\mathcal{U}}
\newcommand{\cX}{\mathcal{X}}
\newcommand{\cG}{\mathcal{G}}
\newcommand{\cL}{\mathcal{L}}
\newcommand{\cB}{\mathcal{B}}
\newcommand{\cZ}{\mathcal{Z}}
\newcommand{\cD}{\mathcal{D}}
\newcommand{\ip}[2]{\langle #1, #2 \rangle}
\newcommand{\norm}[1]{\|#1\|}
\newcommand{\abs}[1]{|#1|}
\newcommand{\btheta}{\bm{\theta}}
\newcommand{\bmu}{\bm{\mu}}
\newcommand{\bu}{\bm{u}}
\newcommand{\bw}{\bm{w}}
\newcommand{\bx}{\bm{x}}
\newcommand{\bxi}{\bm{\xi}}
\newcommand{\bL}{\bm{L}}
\newcommand{\bg}{\bm{g}}
\newcommand{\br}{\bm{r}}

\DeclareMathOperator{\argmin}{argmin}
\DeclareMathOperator{\argmax}{argmax}
\DeclareMathOperator{\ri}{ri}
\DeclareMathOperator{\tr}{tr}
\DeclareMathOperator{\diag}{diag}

% --- Variational analysis helpers (Clarke generalized gradients) ---
\DeclareMathOperator{\conv}{conv}
\DeclareMathOperator{\cl}{cl}
\newcommand{\clsubd}{\partial^{\mathrm{C}}} % Clarke subdifferential
\newcommand{\ipF}[2]{\langle #1, #2 \rangle_F} % Frobenius inner product
\newcommand{\cM}{\mathcal{M}} % Dual multiplier set

%===============================================================================
% Theorem Environments
%===============================================================================
\newtheorem{theorem}{Theorem}
\newtheorem{lemma}[theorem]{Lemma}
\newtheorem{proposition}[theorem]{Proposition}
\newtheorem{corollary}[theorem]{Corollary}

\theoremstyle{definition}
\newtheorem{definition}{Definition}
\newtheorem{assumption}{Assumption}
\newtheorem{remark}{Remark}

%===============================================================================
% Hyphenation
%===============================================================================
\hyphenation{op-ti-mi-za-tion un-cer-tain-ty}

%===============================================================================
% Document
%===============================================================================
\begin{document}

\title{Differentiable Robust Optimization for Learning Correlated Uncertainty Sets in Operating Reserve Markets}

\author{%
\thanks{Manuscript received XXX; revised XXX. This work was supported in part by XXX.}%
}

\markboth{IEEE Transactions on Power Systems}%
{Differentiable Robust Optimization for Learning Correlated Uncertainty Sets}

\maketitle

%===============================================================================
% Abstract
%===============================================================================
\begin{abstract}
Electricity markets rely on operating reserve procurement to manage uncertainty arising from renewable generation. In current practice, reserve requirements are typically determined using simplified uncertainty models that neglect the strong spatial correlations induced by weather. This paper develops a differentiable, learning-based robust optimization framework in which the geometry of ellipsoidal uncertainty sets is learned end-to-end from data. The framework profiles the bilevel problem---minimizing robust dispatch cost subject to coverage constraints---into a single-level objective, and derives the true gradient including a quantile-sensitivity correction term. For unknown distributions, a three-way data split (train/tune/calibrate) enables gradient approximation while preserving distribution-free coverage guarantees via split conformal prediction. Two structural regimes are identified: zonal reserve formulations admit closed-form gradients without solving optimization problems, while network-constrained formulations require envelope gradients from dual information. The framework is instantiated on the IEEE~118-bus system, demonstrating how learned correlation structure affects reserve procurement efficiency.
\end{abstract}

\begin{IEEEkeywords}
Robust optimization, uncertainty sets, operating reserves, machine learning, conformal prediction, electricity markets.
\end{IEEEkeywords}

%===============================================================================
% I. Introduction
%===============================================================================
\section{Introduction}
\label{sec:introduction}

\IEEEPARstart{E}{lectricity} markets rely on operating reserve procurement and pricing mechanisms to manage uncertainty arising from high penetrations of wind and solar generation. Renewable forecast errors are not only substantial in magnitude but also exhibit strong spatial and temporal correlations due to common weather patterns~\cite{Hodge2011,Pinson2013}. When such correlations are ignored, reserve capacity requirements may be inefficiently allocated across the power network, leading to excessive procurement costs, distorted price signals, and reduced overall market efficiency.

Robust optimization provides a principled framework for incorporating uncertainty into power system operations and market clearing~\cite{BenTal2009,Bertsimas2011}. In security-constrained economic dispatch (SCED) and unit commitment, uncertain parameters---such as net load forecast errors---are assumed to lie within a prescribed \emph{uncertainty set}, and the dispatch solution must remain feasible for all realizations within this set. The geometry of the uncertainty set fundamentally determines the conservatism and cost of the resulting solution~\cite{Bertsimas2013}.

In most practical implementations, uncertainty sets are specified with fixed geometric structure (e.g., polyhedral boxes, spherical balls, or ellipsoids with prescribed covariance), and only their overall size is tuned to meet reliability criteria~\cite{Jiang2012,Zeng2013}. While this approach supports tractable market clearing, it limits the ability of reserve and dispatch mechanisms to adapt to spatially and temporally varying correlation patterns observed in renewable forecast errors.

\subsection{Literature Review}

Robust optimization for power system operations has been extensively studied. Bertsimas et al.~\cite{Bertsimas2013} developed adaptive robust optimization for security-constrained unit commitment with polyhedral uncertainty sets. Jiang et al.~\cite{Jiang2012} proposed two-stage robust optimization for unit commitment with uncertain wind power. These approaches typically assume fixed uncertainty set geometry determined a priori.

Data-driven approaches to uncertainty set construction have gained attention. Bertsimas et al.~\cite{Bertsimas2018} proposed data-driven robust optimization that constructs uncertainty sets from historical data using hypothesis testing. Esfahani and Kuhn~\cite{Esfahani2018} developed distributionally robust optimization with Wasserstein ambiguity sets. However, these methods do not directly optimize the uncertainty set geometry for downstream decision-making costs.

Recent work on differentiable optimization has enabled end-to-end learning through optimization layers~\cite{Amos2017,Agrawal2019}. These techniques differentiate through the solution of convex programs, enabling gradient-based learning of problem parameters. However, differentiating through optimization solvers can be computationally expensive and numerically challenging.

Conformal prediction provides distribution-free uncertainty quantification with finite-sample coverage guarantees~\cite{Vovk2005}. Romano et al.~\cite{Romano2019} developed conformalized quantile regression, and Tibshirani et al.~\cite{Tibshirani2019} extended conformal prediction to handle covariate shift. These methods have not been integrated with robust optimization for power systems applications.

\subsection{Contributions}

This paper proposes a learning-based robust optimization framework for operating reserve procurement in which the geometry of uncertainty sets is learned directly from data. The key contributions are:

\begin{enumerate}
\item \textbf{General Differentiable Robust Formulation:} A bilevel optimization framework using gauge-based uncertainty sets $\cU_{\btheta,\rho} = \{\bu : s_{\btheta}(\bu) \leq \rho\}$, where shape $\btheta$ and size $\rho$ are co-optimized.

\item \textbf{True Profiled Gradient:} Derivation of the gradient of the profiled objective $J_P(\btheta) = V(\btheta, \rho_P(\btheta))$, which includes both a direct shape effect and a quantile-sensitivity correction term that depends on the unknown distribution~$P$.

\item \textbf{Tuning-Based Gradient Approximation:} A three-way data split (train/tune/calibrate) with smoothed CDF estimation to approximate the profiled gradient while preserving distribution-free conformal coverage guarantees.

\item \textbf{Envelope-Based Gradients:} Derivation of gradients without differentiating through optimization solvers, using optimal dual multipliers from the robust problem.

\item \textbf{Two Structural Regimes:} Identification of zonal reserve formulations that admit exact decomposition with closed-form gradients (solver-free) and network-constrained formulations that require envelope gradients.

\item \textbf{Power Systems Application:} Instantiation on the IEEE~118-bus system with security-constrained economic dispatch, demonstrating how correlation structure affects reserve requirements.
\end{enumerate}

\subsection{Paper Organization}

The remainder of this paper is organized as follows. Section~\ref{sec:formulation} presents the general problem formulation with gauge-based uncertainty sets. Section~\ref{sec:methodology} develops the methodology, including envelope gradients, profiled gradients, and conformal calibration. Section~\ref{sec:case-study} instantiates the framework for the IEEE~118-bus system. Section~\ref{sec:algorithm} presents the complete algorithm. Section~\ref{sec:discussion} discusses practical implications and correlation effects. Section~\ref{sec:conclusion} concludes.

%===============================================================================
% II. Problem Formulation
%===============================================================================
\section{Problem Formulation}
\label{sec:formulation}

This section presents the general formulation for robust optimization with parameterized uncertainty sets. The framework separates uncertainty set \emph{shape} (controlled by parameter $\btheta$) from \emph{size} (controlled by radius $\rho$), enabling independent optimization of geometry and coverage.

\subsection{Gauge-Based Uncertainty Sets}

Fix an uncertainty dimension $d \geq 1$ and a parameter set $\Theta \subseteq \R^p$ (shape space). Uncertainty sets are defined via the gauge (Minkowski functional) of a unit set.

\begin{definition}[Parameterized Uncertainty Set]
For shape $\btheta \in \Theta$ and radius $\rho > 0$, the uncertainty set is
\begin{equation}\label{eq:uset}
\cU_{\btheta,\rho} := \{\bu \in \R^d : s_{\btheta}(\bu) \leq \rho\}
\end{equation}
where $s_{\btheta}: \R^d \to [0,\infty)$ is the gauge function of the unit set $\cU_{\btheta,1}$.
\end{definition}

\begin{assumption}[Gauge Regularity]\label{asm:gauge}
For each $\btheta \in \Theta$, the unit set $\cU_{\btheta,1}$ is nonempty, convex, compact, and contains $\bm{0}$ in its interior. The gauge function satisfies
\begin{equation}
s_{\btheta}(\bu) = \inf\{t > 0 : \bu \in t \cdot \cU_{\btheta,1}\}.
\end{equation}
Consequently, $\cU_{\btheta,\rho} = \rho \cdot \cU_{\btheta,1}$ for all $\rho > 0$.
\end{assumption}

\begin{definition}[Support Function]
For a nonempty set $C \subseteq \R^d$, the support function is
\begin{equation}
\sigma_C(\bw) := \sup_{\bu \in C} \ip{\bw}{\bu} \in (-\infty, +\infty].
\end{equation}
\end{definition}

The support function is the key object for reformulating robust constraints. The following scaling property is fundamental.

\begin{lemma}[Support Function Scaling]\label{lem:support-scaling}
If $C \subseteq \R^d$ is nonempty and $\rho > 0$, then
\begin{equation}\label{eq:sigma-scale}
\sigma_{\rho C}(\bw) = \rho \cdot \sigma_C(\bw) \quad \forall \bw \in \R^d.
\end{equation}
\end{lemma}

\begin{proof}
By definition, $\sigma_{\rho C}(\bw) = \sup_{\bu \in \rho C} \ip{\bw}{\bu} = \sup_{\bm{v} \in C} \ip{\bw}{\rho \bm{v}} = \rho \sup_{\bm{v} \in C} \ip{\bw}{\bm{v}} = \rho \sigma_C(\bw)$.
\end{proof}

Under Assumption~\ref{asm:gauge}, Lemma~\ref{lem:support-scaling} yields
\begin{equation}
\sigma_{\cU_{\btheta,\rho}}(\bw) = \rho \cdot \sigma_{\cU_{\btheta,1}}(\bw).
\end{equation}

\subsection{Robust Optimization Value Function}

Let $\cX \subseteq \R^{n_x}$ be a closed convex decision set, $f: \R^{n_x} \to \R \cup \{+\infty\}$ be a proper closed convex objective, and $a_j: \R^{n_x} \to \R \cup \{+\infty\}$ be proper closed convex constraint functions. Fix exposure vectors $\bw_1, \ldots, \bw_m \in \R^d$ and scalars $b_1, \ldots, b_m \in \R$.

Given shape $\btheta$ and radius $\rho$, the (extended) robust value function is
\begin{equation}\label{eq:robust}
\begin{aligned}
V(\btheta,\rho) := \inf_{\bx \in \cX} \quad & f(\bx) \\
\text{s.t.} \quad & a_j(\bx) + \sigma_{\cU_{\btheta,\rho}}(\bw_j) \leq b_j, \quad j = 1,\ldots,m,
\end{aligned}
\end{equation}
with the convention $V(\btheta,\rho)=+\infty$ if the feasible set is empty.
The constraints require that $a_j(\bx) + \ip{\bw_j}{\bu} \leq b_j$ holds for all $\bu \in \cU_{\btheta,\rho}$, which reduces to (\ref{eq:robust}) via the support function.

\begin{assumption}[Local Slater and Dual Attainment]\label{asm:slater}
Fix $(\btheta,\rho)\in\Theta\times(0,\infty)$ such that $V(\btheta,\rho)<+\infty$.
Assume: (i) $\cX$ is nonempty, closed, convex; (ii) $f$ and all $a_j$ are proper closed convex; and (iii) there exists $\bar{\bx} \in \ri(\cX)$ such that
\begin{equation}\label{eq:slater-strict}
a_j(\bar{\bx}) + \sigma_{\cU_{\btheta,\rho}}(\bw_j) < b_j, \quad j = 1,\ldots,m.
\end{equation}
Moreover, assume that the Lagrange dual problem (defined below) attains an optimal solution, i.e., the set of optimal dual multipliers $\cM^*(\btheta,\rho)$ is nonempty.
\end{assumption}

Define the constraint functions
\[
g_j(\bx;\btheta,\rho) := a_j(\bx) + \sigma_{\cU_{\btheta,\rho}}(\bw_j) - b_j.
\]
The Lagrangian with multipliers $\bmu\in\R_+^m$ is
\[
\cL(\bx,\bmu;\btheta,\rho)= f(\bx) + I_{\cX}(\bx) + \sum_{j=1}^m \mu_j g_j(\bx;\btheta,\rho),
\]
where $I_{\cX}(\bx) = 0$ if $\bx \in \cX$ and $+\infty$ otherwise. The dual function is
\[
q(\bmu;\btheta,\rho) := \inf_{\bx\in\R^{n_x}} \cL(\bx,\bmu;\btheta,\rho).
\]
The dual problem is $\max_{\bmu\ge 0} q(\bmu;\btheta,\rho)$, and the argmax set is denoted $\cM^*(\btheta,\rho)$.

\begin{lemma}[Strong Duality at Slater Points]\label{lem:strong-duality}
Under Assumption~\ref{asm:slater}, strong duality holds at $(\btheta,\rho)$:
\[
V(\btheta,\rho)=\max_{\bmu\ge 0} q(\bmu;\btheta,\rho),
\]
and the maximum is attained, i.e., $\cM^*(\btheta,\rho)\neq\emptyset$.
\end{lemma}

\subsection{The Bilevel Learning Problem}

Let $U \sim P$ be a random uncertainty realization in $\R^d$. Define the gauge score random variable $S_{\btheta} := s_{\btheta}(U)$ and its cumulative distribution function (CDF) under $P$:
\begin{equation}
F_{\btheta}(r) := \Prob(S_{\btheta} \leq r), \quad r \in \R.
\end{equation}
Fix a target coverage level $\tau \in (0,1)$ (e.g., $\tau = 0.95$).

The learning problem seeks to minimize the robust value function subject to a probabilistic coverage constraint:
\begin{equation}\label{eq:bilevel}
\boxed{
\min_{\btheta \in \Theta,\, \rho > 0} V(\btheta,\rho) \quad \text{s.t.} \quad \Prob(U \in \cU_{\btheta,\rho}) \geq \tau
}
\end{equation}

Since $U \in \cU_{\btheta,\rho}$ if and only if $s_{\btheta}(U) \leq \rho$, the coverage constraint is equivalent to $F_{\btheta}(\rho) \geq \tau$.

Define the $\tau$-quantile radius map:
\begin{equation}\label{eq:rhoP}
\rho_P(\btheta) := \inf\{r > 0 : F_{\btheta}(r) \geq \tau\}.
\end{equation}

\begin{proposition}[Radius Profiling]\label{prop:profiling}
Under Assumptions~\ref{asm:gauge} and~\ref{asm:slater}, fix $\btheta \in \Theta$ and suppose $\rho_P(\btheta) < \infty$. Then:
\begin{enumerate}
\item The map $\rho \mapsto V(\btheta, \rho)$ is nondecreasing on $(0,\infty)$.
\item The smallest feasible radius for the coverage constraint is $\rho = \rho_P(\btheta)$.
\item The bilevel problem~(\ref{eq:bilevel}) reduces to the profiled problem
\begin{equation}\label{eq:profiled}
\min_{\btheta \in \Theta} J_P(\btheta) \quad \text{where} \quad J_P(\btheta) := V(\btheta, \rho_P(\btheta)).
\end{equation}
\end{enumerate}
\end{proposition}

\begin{proof}
(i) Fix $\btheta$ and let $0 < \rho_1 \leq \rho_2$. Since $\cU_{\btheta,\rho_1} \subseteq \cU_{\btheta,\rho_2}$, the support function satisfies $\sigma_{\cU_{\btheta,\rho_1}}(\bw_j) \leq \sigma_{\cU_{\btheta,\rho_2}}(\bw_j)$. Thus the feasible region of~(\ref{eq:robust}) shrinks as $\rho$ increases, and $V(\btheta,\rho_1) \leq V(\btheta,\rho_2)$.

(ii) By definition of $\rho_P(\btheta)$, the coverage constraint $F_{\btheta}(\rho) \geq \tau$ holds if and only if $\rho \geq \rho_P(\btheta)$.

(iii) For any feasible $(\btheta, \rho)$ in~(\ref{eq:bilevel}), we have $\rho \geq \rho_P(\btheta)$ by (ii). By (i), $V(\btheta,\rho) \geq V(\btheta,\rho_P(\btheta)) = J_P(\btheta)$. Thus the optimal choice for each $\btheta$ is $\rho = \rho_P(\btheta)$.
\end{proof}

%===============================================================================
% III. Methodology
%===============================================================================
\section{Methodology}
\label{sec:methodology}

This section develops the complete methodology for optimizing the profiled objective~(\ref{eq:profiled}). The key challenge is that both the robust value function $V(\btheta,\rho)$ and the quantile map $\rho_P(\btheta)$ depend on the shape parameter $\btheta$, and the quantile map depends on the unknown distribution~$P$.

\subsection{Envelope Gradients of $V(\btheta,\rho)$}

A central feature of the proposed approach is that gradients of the robust value function are obtained without differentiating through the optimization solver. Instead, the envelope theorem provides gradients in terms of optimal dual multipliers. When dual multipliers may not be unique, Clarke's generalized subdifferential provides the appropriate notion.

\begin{definition}[Clarke Subdifferential]\label{def:clarke}
Let $h:\R^p\to\R$ be locally Lipschitz. The Clarke subdifferential of $h$ at $x$ is
\[
\clsubd h(x)
:=\conv\Big\{\lim_{k\to\infty}\nabla h(x_k)\ :\ x_k\to x,\ h \text{ differentiable at }x_k\Big\}.
\]
Any $g\in\clsubd h(x)$ is called a Clarke subgradient of $h$ at $x$.
\end{definition}

Define $\sigma_j(\btheta,\rho) := \sigma_{\cU_{\btheta,\rho}}(\bw_j)$ for shorthand.

\begin{assumption}[Support Function Smoothness]\label{asm:sigma-diff}
For each $j \in \{1,\ldots,m\}$ and $\rho > 0$, the map $\btheta \mapsto \sigma_j(\btheta,\rho)$ is continuously differentiable on $\Theta$, and $\nabla_{\btheta}\sigma_j(\btheta,\rho)$ is locally bounded on $\Theta$.
\end{assumption}

\begin{theorem}[Envelope Formula via Dual Multipliers (Clarke Form)]\label{thm:envelope-theta}
Under Assumptions~\ref{asm:slater} and~\ref{asm:sigma-diff}, fix $(\btheta,\rho)$ with $V(\btheta,\rho)<+\infty$.
Let $\cM^*(\btheta,\rho)$ denote the set of optimal dual multipliers.

Then for any $\bmu^*\in\cM^*(\btheta,\rho)$,
\begin{equation}\label{eq:env-theta}
\bg_{\btheta}(\btheta,\rho;\bmu^*) := \sum_{j=1}^m \mu_j^* \nabla_{\btheta} \sigma_{\cU_{\btheta,\rho}}(\bw_j)
\end{equation}
is a Clarke subgradient of $\btheta\mapsto V(\btheta,\rho)$ at $\btheta$, i.e.,
\[
\bg_{\btheta}(\btheta,\rho;\bmu^*)\in \clsubd_{\btheta} V(\btheta,\rho).
\]

If, in addition, the dual optimizer is unique at $(\btheta,\rho)$ (i.e., $\cM^*(\btheta,\rho)=\{\bmu^*(\btheta,\rho)\}$),
then $V(\cdot,\rho)$ is differentiable at $\btheta$ and
\[
\nabla_{\btheta} V(\btheta,\rho)=\sum_{j=1}^m \mu_j^*(\btheta,\rho)\,\nabla_{\btheta}\sigma_{\cU_{\btheta,\rho}}(\bw_j).
\]
\end{theorem}

\begin{proof}
Fix $\bmu^*\in\cM^*(\btheta,\rho)$.
For any $\btheta'$,
\[
V(\btheta',\rho)=\max_{\bmu\ge 0} q(\bmu;\btheta',\rho)\ \ge\ q(\bmu^*;\btheta',\rho).
\]
Because $\sigma_{\cU_{\btheta,\rho}}(\bw_j)$ does not depend on $\bx$, the dual function separates as
\[
q(\bmu;\btheta,\rho)=\tilde q(\bmu)+\sum_{j=1}^m \mu_j\big(\sigma_{\cU_{\btheta,\rho}}(\bw_j)-b_j\big),
\]
where $\tilde q(\bmu) = \inf_{\bx}\{f(\bx) + I_{\cX}(\bx) + \sum_j \mu_j a_j(\bx)\}$ is independent of $\btheta$.
Thus
\[
q(\bmu^*;\btheta',\rho)-q(\bmu^*;\btheta,\rho)
=\sum_{j=1}^m \mu_j^*\Big(\sigma_{\cU_{\btheta',\rho}}(\bw_j)-\sigma_{\cU_{\btheta,\rho}}(\bw_j)\Big).
\]
Using $q(\bmu^*;\btheta,\rho)=V(\btheta,\rho)$ and the differentiability of each
$\btheta\mapsto\sigma_{\cU_{\btheta,\rho}}(\bw_j)$ yields
\[
V(\btheta',\rho)-V(\btheta,\rho)
\ge \ip{\bg_{\btheta}(\btheta,\rho;\bmu^*)}{\btheta'-\btheta} + o(\|\btheta'-\btheta\|),
\]
which implies $\bg_{\btheta}(\btheta,\rho;\bmu^*)\in \clsubd_{\btheta}V(\btheta,\rho)$ by Definition~\ref{def:clarke}.
The differentiable case under uniqueness is standard: the Clarke subdifferential reduces to a singleton, equal to the gradient.
\end{proof}

\begin{theorem}[Envelope Formula for Radius (Clarke Form)]\label{thm:envelope-rho}
Under Assumptions~\ref{asm:gauge} and~\ref{asm:slater}, fix $(\btheta,\rho)$ with $V(\btheta,\rho)<+\infty$,
and let $\cM^*(\btheta,\rho)$ be the optimal dual set.

For any $\bmu^*\in\cM^*(\btheta,\rho)$,
\begin{equation}\label{eq:env-rho}
g_{\rho}(\btheta,\rho;\bmu^*)
:=\sum_{j=1}^m \mu_j^*\,\sigma_{\cU_{\btheta,1}}(\bw_j)
\end{equation}
is a Clarke subgradient of $\rho\mapsto V(\btheta,\rho)$ at $\rho$, i.e.,
$g_\rho(\btheta,\rho;\bmu^*)\in \clsubd_{\rho}V(\btheta,\rho)$.

If the dual optimizer is unique at $(\btheta,\rho)$, then $V(\btheta,\cdot)$ is differentiable at $\rho$ and
\[
\frac{\partial}{\partial\rho}V(\btheta,\rho)
=\sum_{j=1}^m \mu_j^*(\btheta,\rho)\,\sigma_{\cU_{\btheta,1}}(\bw_j).
\]
\end{theorem}

\begin{proof}
By Assumption~\ref{asm:gauge}, $\cU_{\btheta,\rho}=\rho\,\cU_{\btheta,1}$ and hence
$\sigma_{\cU_{\btheta,\rho}}(\bw_j)=\rho\,\sigma_{\cU_{\btheta,1}}(\bw_j)$.
Therefore $q(\bmu;\btheta,\rho)$ is affine in $\rho$:
\[
q(\bmu;\btheta,\rho)
=\tilde q(\bmu)-\sum_{j=1}^m \mu_j b_j
+\rho\sum_{j=1}^m \mu_j\sigma_{\cU_{\btheta,1}}(\bw_j).
\]
The same argument as in Theorem~\ref{thm:envelope-theta} yields the Clarke subgradient statement.
\end{proof}

\begin{remark}[No Differentiation Through Solver]
Theorems~\ref{thm:envelope-theta} and~\ref{thm:envelope-rho} require only the optimal dual multipliers $\bmu^*$ from solving~(\ref{eq:robust}), not the derivative of the solution mapping. This avoids the computational and numerical challenges of differentiating through optimization solvers~\cite{Amos2017,Agrawal2019}.
\end{remark}

\subsection{The True Profiled Gradient (Oracle Case)}

When the distribution $P$ is known, the true gradient of $J_P(\btheta) = V(\btheta, \rho_P(\btheta))$ follows from the chain rule, incorporating both the direct effect of $\btheta$ on $V$ and the indirect effect through $\rho_P(\btheta)$.

\begin{assumption}[Quantile Regularity]\label{asm:quantile-reg}
Fix $\btheta \in \Theta$ and let $\rho_P(\btheta)$ be defined as in~(\ref{eq:rhoP}). There exists an open interval $I$ containing $\rho_P(\btheta)$ such that: (i) $F_{\btheta}$ is continuously differentiable on $I$ with density $f_{\btheta}(r) := \partial_r F_{\btheta}(r)$; (ii) $f_{\btheta}(\rho_P(\btheta)) > 0$; and (iii) for $\btheta'$ in a neighborhood of $\btheta$, the map $(\btheta',r) \mapsto F_{\btheta'}(r)$ is continuously differentiable.
\end{assumption}

\begin{lemma}[Implicit Derivative of Quantile Map]\label{lem:quantile-deriv}
Under Assumption~\ref{asm:quantile-reg}, $\rho_P$ is differentiable at $\btheta$ with
\begin{equation}\label{eq:quantile-ift}
\nabla_{\btheta} \rho_P(\btheta) = -\frac{\nabla_{\btheta} F_{\btheta}(\rho_P(\btheta))}{f_{\btheta}(\rho_P(\btheta))}.
\end{equation}
\end{lemma}

\begin{proof}
Define $G(\btheta,r) := F_{\btheta}(r) - \tau$. By definition, $G(\btheta, \rho_P(\btheta)) = 0$ and $\partial_r G = f_{\btheta}(\rho_P(\btheta)) > 0$. The implicit function theorem yields~(\ref{eq:quantile-ift}).
\end{proof}

\begin{theorem}[True Oracle Gradient]\label{thm:true-oracle-grad}
Assume Assumptions~\ref{asm:gauge}, \ref{asm:slater}, \ref{asm:sigma-diff}, and~\ref{asm:quantile-reg} hold, and that $V(\cdot,\rho)$ is differentiable at $\btheta$ for $\rho = \rho_P(\btheta)$ and $V(\btheta,\cdot)$ is differentiable at $\rho_P(\btheta)$. Then $J_P$ is differentiable at $\btheta$ with
\begin{equation}\label{eq:true-grad}
\boxed{
\nabla_{\btheta} J_P(\btheta) = \nabla_{\btheta} V(\btheta, \rho_P(\btheta)) + \frac{\partial V}{\partial\rho}(\btheta, \rho_P(\btheta)) \cdot \nabla_{\btheta} \rho_P(\btheta)
}
\end{equation}
\end{theorem}

\begin{proof}
Apply the chain rule to $J_P(\btheta) = V(\btheta, \rho_P(\btheta))$.
\end{proof}

Let $\bmu^* = \bmu^*(\btheta, \rho_P(\btheta))$ denote the optimal dual multipliers at the profiled point. Substituting the envelope gradients from Theorems~\ref{thm:envelope-theta}--\ref{thm:envelope-rho} and Lemma~\ref{lem:quantile-deriv}:
\begin{equation}\label{eq:true-grad-expanded}
\boxed{
\begin{aligned}
\nabla_{\btheta} J_P(\btheta) &= \underbrace{\sum_{j=1}^m \mu_j^* \nabla_{\btheta} \sigma_{\cU_{\btheta,\rho_P(\btheta)}}(\bw_j)}_{\text{direct shape effect}} \\
&\quad + \underbrace{\left(\sum_{j=1}^m \mu_j^* \sigma_{\cU_{\btheta,1}}(\bw_j)\right) \left(-\frac{\nabla_{\btheta} F_{\btheta}(\rho_P(\btheta))}{f_{\btheta}(\rho_P(\btheta))}\right)}_{\text{quantile-sensitivity correction}}
\end{aligned}
}
\end{equation}

\begin{remark}[Where $P$ Enters the Gradient]
The second term in~(\ref{eq:true-grad-expanded}) is the \emph{quantile-sensitivity correction}. It depends on $P$ through $\nabla_{\btheta} F_{\btheta}(\rho_P(\btheta))$ and $f_{\btheta}(\rho_P(\btheta))$. If $P$ is unknown, this term cannot be computed exactly. This is why naive ``optimize $\btheta$ then calibrate $\rho$'' does not literally optimize the oracle profiled objective $J_P$.
\end{remark}

\subsection{Distribution-Free Conformal Calibration}

When $P$ is unknown, split conformal prediction provides distribution-free coverage guarantees using an independent calibration sample.

Let $\{U_i\}_{i=1}^{n_{\text{cal}}}$ be a calibration sample.

\begin{assumption}[Calibration Exchangeability]\label{asm:exchange}
Conditional on the learned shape $\hat{\btheta}$, the calibration sample
$(U_1,\ldots,U_{n_{\mathrm{cal}}})$ and the future realization $U_{\mathrm{new}}$ are exchangeable.
In particular, this holds if $(U_i)_{i=1}^{n_{\mathrm{cal}}}$ and $U_{\mathrm{new}}$ are i.i.d.\ from the same distribution
and $\hat{\btheta}$ is independent of the calibration sample.
\end{assumption}

\begin{theorem}[Split Conformal Coverage]\label{thm:conformal}
Under Assumption~\ref{asm:exchange}, fix any (possibly random) shape $\hat{\btheta}$ that is independent of the calibration sample $(U_i)_{i=1}^{n_{\text{cal}}}$. Define calibration scores $S_i = s_{\hat{\btheta}}(U_i)$ and let $S_{(1)} \leq \cdots \leq S_{(n_{\text{cal}})}$ be their order statistics. Set $S_{(n_{\text{cal}}+1)} := +\infty$ and define
\begin{equation}\label{eq:conformal-rho}
k := \lceil (n_{\text{cal}}+1)\tau \rceil, \quad \hat{\rho}_\tau := S_{(k)}.
\end{equation}
Then
\begin{equation}
\Prob(s_{\hat{\btheta}}(U_{\text{new}}) \leq \hat{\rho}_\tau) \geq \tau
\end{equation}
equivalently, $\Prob(U_{\text{new}} \in \cU_{\hat{\btheta}, \hat{\rho}_\tau}) \geq \tau$.
\end{theorem}

\begin{proof}
Consider the $(n_{\text{cal}}+1)$ scores $S_1, \ldots, S_{n_{\text{cal}}}, S_{\text{new}}$ where $S_{\text{new}} := s_{\hat{\btheta}}(U_{\text{new}})$. By exchangeability and independence of $\hat{\btheta}$ from the calibration sample, these scores are exchangeable. Thus the rank of $S_{\text{new}}$ among all $(n_{\text{cal}}+1)$ scores is uniform on $\{1, \ldots, n_{\text{cal}}+1\}$. The event $S_{\text{new}} \leq S_{(k)}$ occurs when $\text{rank}(S_{\text{new}}) \leq k$, hence
\begin{equation}
\Prob(S_{\text{new}} \leq S_{(k)}) = \frac{k}{n_{\text{cal}}+1} \geq \tau
\end{equation}
since $k = \lceil (n_{\text{cal}}+1)\tau \rceil$.
\end{proof}

\begin{remark}[Independence Requirement]
The coverage guarantee requires only that $\hat{\btheta}$ be independent of $\cD_{\text{cal}}$. The shape may depend on separate training and tuning data; the conformal guarantee holds regardless of how $\hat{\btheta}$ was learned.
\end{remark}

\begin{remark}[Time-Series Data and Validity of Exchangeability]
The finite-sample coverage guarantee in Theorem~\ref{thm:conformal} is exact under Assumption~\ref{asm:exchange}.
If the data arise from a temporally dependent time series, exchangeability may fail and the theorem does not apply as stated.
To retain a rigorous finite-sample guarantee, one must construct calibration examples that satisfy exchangeability
(e.g., by working with independent or nonoverlapping blocks that are assumed exchangeable, or by using data drawn i.i.d.\ across operating intervals).
Absent such a construction or assumption, conformal coverage should be treated as approximate rather than distribution-free.
\end{remark}

\subsection{Tuning-Based Gradient Approximation}

To approximate the quantile-sensitivity correction when $P$ is unknown, the data is partitioned into three disjoint sets:
\begin{itemize}
\item $\cD_{\text{train}}$: optimize shape parameters via gradient descent
\item $\cD_{\text{tune}}$: approximate quantile sensitivity $\nabla_{\btheta} \rho_P(\btheta)$
\item $\cD_{\text{cal}}$: final conformal radius calibration
\end{itemize}

The main obstacle is that $F_{\btheta}(r) = \E[\mathbf{1}\{s_{\btheta}(U) \leq r\}]$ involves a non-differentiable indicator. A standard smooth approximation is employed.

\begin{assumption}[Smoothing Function]\label{asm:smoothing}
Let $\Phi: \R \to [0,1]$ be continuously differentiable, nondecreasing, with $\lim_{t \to -\infty} \Phi(t) = 0$ and $\lim_{t \to +\infty} \Phi(t) = 1$. Let $\varphi(t) := \Phi'(t)$ be bounded.
\end{assumption}

For bandwidth $\varepsilon > 0$, define the smoothed CDF:
\begin{equation}\label{eq:smoothed-cdf}
F_{\btheta,\varepsilon}(r) := \E\left[\Phi\left(\frac{r - s_{\btheta}(U)}{\varepsilon}\right)\right].
\end{equation}

\begin{assumption}[Score Differentiability]\label{asm:score-diff}
For each $\bu \in \R^d$, the map $\btheta \mapsto s_{\btheta}(\bu)$ is continuously differentiable. There exists an integrable function $M(U)$ such that $\norm{\nabla_{\btheta} s_{\btheta}(U)} \leq M(U)$ almost surely.
\end{assumption}

\begin{lemma}[Smoothed CDF Derivatives]\label{lem:smoothed-derivatives}
Under Assumptions~\ref{asm:smoothing} and~\ref{asm:score-diff}, for each $\varepsilon > 0$:
\begin{align}
\partial_r F_{\btheta,\varepsilon}(r) &= \frac{1}{\varepsilon} \E\left[\varphi\left(\frac{r - s_{\btheta}(U)}{\varepsilon}\right)\right] \label{eq:dFr} \\
\nabla_{\btheta} F_{\btheta,\varepsilon}(r) &= -\frac{1}{\varepsilon} \E\left[\varphi\left(\frac{r - s_{\btheta}(U)}{\varepsilon}\right) \nabla_{\btheta} s_{\btheta}(U)\right] \label{eq:dFtheta}
\end{align}
\end{lemma}

\begin{proof}
Differentiation under the expectation is justified by dominated convergence using boundedness of $\varphi$ and integrability of $\nabla_{\btheta} s_{\btheta}(U)$.
\end{proof}

Define the smoothed $\tau$-quantile radius:
\begin{equation}
\rho_{P,\varepsilon}(\btheta) := \inf\{r : F_{\btheta,\varepsilon}(r) \geq \tau\}.
\end{equation}

\begin{lemma}[Smoothed Quantile Sensitivity]\label{lem:smoothed-quantile-deriv}
Under Assumptions~\ref{asm:smoothing}, \ref{asm:score-diff}, and assuming $\partial_r F_{\btheta,\varepsilon}(\rho_{P,\varepsilon}(\btheta)) > 0$:
\begin{equation}\label{eq:rhoeps-deriv}
\nabla_{\btheta} \rho_{P,\varepsilon}(\btheta) = \frac{\E\left[\varphi\left(\frac{\rho_{P,\varepsilon}(\btheta) - s_{\btheta}(U)}{\varepsilon}\right) \nabla_{\btheta} s_{\btheta}(U)\right]}{\E\left[\varphi\left(\frac{\rho_{P,\varepsilon}(\btheta) - s_{\btheta}(U)}{\varepsilon}\right)\right]}
\end{equation}
\end{lemma}

\begin{proof}
Apply the implicit function theorem to $G(\btheta,r) = F_{\btheta,\varepsilon}(r) - \tau$ and substitute Lemma~\ref{lem:smoothed-derivatives}. The factor $1/\varepsilon$ cancels in the ratio.
\end{proof}

\begin{remark}[Interpretation]
The formula~(\ref{eq:rhoeps-deriv}) is a weighted average of $\nabla_{\btheta} s_{\btheta}(U)$ with weights $\varphi((\rho_{P,\varepsilon} - s_{\btheta}(U))/\varepsilon)$ concentrated on samples near the quantile boundary. This is a soft analogue of the boundary conditional expectation $\E[\nabla_{\btheta} s_{\btheta}(U) \mid s_{\btheta}(U) = \rho_P(\btheta)]$.
\end{remark}

Given tuning data $\cD_{\text{tune}} = \{U_i\}_{i=1}^{n_{\text{tune}}}$, define empirical scores $S_i(\btheta) = s_{\btheta}(U_i)$ and the empirical smoothed quantile:
\begin{equation}
\hat{\rho}_\varepsilon(\btheta) := \inf\left\{r : \frac{1}{n_{\text{tune}}} \sum_{i=1}^{n_{\text{tune}}} \Phi\left(\frac{r - S_i(\btheta)}{\varepsilon}\right) \geq \tau\right\}.
\end{equation}

With weights $\omega_i(\btheta) := \varphi((\hat{\rho}_\varepsilon(\btheta) - S_i(\btheta))/\varepsilon)$, the empirical quantile sensitivity estimator is:
\begin{equation}\label{eq:rhohat-grad}
\widehat{\nabla_{\btheta} \rho_\varepsilon}(\btheta) := \frac{\sum_{i=1}^{n_{\text{tune}}} \omega_i(\btheta) \nabla_{\btheta} s_{\btheta}(U_i)}{\sum_{i=1}^{n_{\text{tune}}} \omega_i(\btheta)}.
\end{equation}

The approximate profiled gradient estimator combines the envelope terms with the tuned quantile sensitivity:
\begin{equation}\label{eq:ghat}
\boxed{
\begin{aligned}
\hat{\bg}_\varepsilon(\btheta) &:= \underbrace{\sum_{j=1}^m \mu_j^*(\btheta, \hat{\rho}_\varepsilon(\btheta)) \nabla_{\btheta} \sigma_{\cU_{\btheta,\hat{\rho}_\varepsilon(\btheta)}}(\bw_j)}_{\text{envelope shape term}} \\
&\quad + \underbrace{\left(\sum_{j=1}^m \mu_j^*(\btheta, \hat{\rho}_\varepsilon(\btheta)) \sigma_{\cU_{\btheta,1}}(\bw_j)\right)}_{\text{envelope size sensitivity}} \cdot \underbrace{\widehat{\nabla_{\btheta} \rho_\varepsilon}(\btheta)}_{\text{tuned quantile sensitivity}}
\end{aligned}
}
\end{equation}

\begin{theorem}[Consistency of the Tuned Smoothed Profiled Gradient]\label{thm:consistency-eps}
Fix $\varepsilon>0$ and a shape $\btheta\in\Theta$.
Assume:

\begin{enumerate}
\item[(A1)] $U_1,\ldots,U_{n_{\mathrm{tune}}}$ are i.i.d.\ from $P$.
\item[(A2)] Assumptions~\ref{asm:smoothing} and~\ref{asm:score-diff} hold, and $\varphi$ is continuous and strictly positive: $\varphi(t)>0$ for all $t\in\R$.
\item[(A3)] The population smoothed CDF $r\mapsto F_{\btheta,\varepsilon}(r)$ is strictly increasing in a neighborhood of
$\rho_{P,\varepsilon}(\btheta)$ and $\partial_r F_{\btheta,\varepsilon}(\rho_{P,\varepsilon}(\btheta))>0$.
\item[(A4)] The robust value function is differentiable at $(\btheta,\rho_{P,\varepsilon}(\btheta))$ with respect to $\btheta$ and $\rho$,
and the dual optimizer is unique in a neighborhood of $(\btheta,\rho_{P,\varepsilon}(\btheta))$ so that
$\bmu^*(\btheta,\rho)$ is well-defined and continuous there.
\item[(A5)] For each $j$, $\nabla_{\btheta}\sigma_{\cU_{\btheta,\rho}}(\bw_j)$ is continuous in $(\btheta,\rho)$ in a neighborhood of
$(\btheta,\rho_{P,\varepsilon}(\btheta))$.
\end{enumerate}

Then, as $n_{\mathrm{tune}}\to\infty$,
\[
\hat{\bg}_\varepsilon(\btheta)\ \xrightarrow{\ \Prob\ }\ \nabla_{\btheta}J_{P,\varepsilon}(\btheta),
\]
where $J_{P,\varepsilon}(\btheta):=V(\btheta,\rho_{P,\varepsilon}(\btheta))$ and $\hat{\bg}_\varepsilon(\btheta)$ is defined in \eqref{eq:ghat}.
\end{theorem}

\begin{proof}
Define the empirical smoothed CDF (tuning sample)
\[
\widehat F_{\btheta,\varepsilon}(r)
:=\frac{1}{n_{\mathrm{tune}}}\sum_{i=1}^{n_{\mathrm{tune}}}
\Phi\!\left(\frac{r-s_{\btheta}(U_i)}{\varepsilon}\right).
\]
Each summand is bounded in $[0,1]$, and for fixed $\btheta$ it is continuous in $r$.
By the law of large numbers, for each fixed $r$,
$\widehat F_{\btheta,\varepsilon}(r)\to F_{\btheta,\varepsilon}(r)$ in probability.

\textbf{Step 1: Consistency of the smoothed quantile.}
Let $\rho:=\rho_{P,\varepsilon}(\btheta)$ and $\hat\rho:=\hat\rho_\varepsilon(\btheta)$.
By (A3), $F_{\btheta,\varepsilon}$ is strictly increasing near $\rho$ with positive derivative, hence it has a locally unique inverse around level $\tau$.
Standard inverse-continuity arguments for monotone functions imply $\hat\rho\to\rho$ in probability
(since $\widehat F_{\btheta,\varepsilon}$ converges pointwise to $F_{\btheta,\varepsilon}$ and both are monotone in $r$).

\textbf{Step 2: Consistency of the tuned quantile sensitivity.}
Define the random weights
\[
\omega_i(\btheta):=\varphi\!\left(\frac{\hat\rho-s_{\btheta}(U_i)}{\varepsilon}\right),
\]
and note $\omega_i(\btheta)>0$ almost surely by (A2), so the denominator in \eqref{eq:rhohat-grad} is nonzero.
Consider the numerator and denominator:
\[
N_n:=\frac{1}{n_{\mathrm{tune}}}\sum_{i=1}^{n_{\mathrm{tune}}}\omega_i(\btheta)\,\nabla_{\btheta}s_{\btheta}(U_i),
\qquad
D_n:=\frac{1}{n_{\mathrm{tune}}}\sum_{i=1}^{n_{\mathrm{tune}}}\omega_i(\btheta).
\]
Because $\hat\rho\to\rho$ in probability, $\varphi$ is continuous and bounded, and $\|\nabla_{\btheta}s_{\btheta}(U)\|$ is integrably dominated (Assumption~\ref{asm:score-diff}),
a Slutsky/continuous-mapping argument yields
\[
N_n \xrightarrow{\Prob} \E\!\left[\varphi\!\left(\frac{\rho-s_{\btheta}(U)}{\varepsilon}\right)\nabla_{\btheta}s_{\btheta}(U)\right],
\quad
D_n \xrightarrow{\Prob} \E\!\left[\varphi\!\left(\frac{\rho-s_{\btheta}(U)}{\varepsilon}\right)\right].
\]
Moreover the limiting denominator is strictly positive by (A3) since
$\partial_rF_{\btheta,\varepsilon}(\rho)=\frac{1}{\varepsilon}\E[\varphi((\rho-s_{\btheta}(U))/\varepsilon)]>0$.
Therefore the ratio converges:
\[
\widehat{\nabla_{\btheta}\rho_\varepsilon}(\btheta)=\frac{N_n}{D_n}
\xrightarrow{\Prob}
\frac{\E\!\left[\varphi\!\left(\frac{\rho-s_{\btheta}(U)}{\varepsilon}\right)\nabla_{\btheta}s_{\btheta}(U)\right]}
{\E\!\left[\varphi\!\left(\frac{\rho-s_{\btheta}(U)}{\varepsilon}\right)\right]}
=\nabla_{\btheta}\rho_{P,\varepsilon}(\btheta),
\]
where the last equality follows from Lemma~\ref{lem:smoothed-quantile-deriv}.

\textbf{Step 3: Continuity of the envelope components.}
By Step 1, $\hat\rho\to\rho$ in probability. By (A4), $\bmu^*(\btheta,\rho)$ is continuous in $\rho$ near $\rho$,
so $\bmu^*(\btheta,\hat\rho)\to\bmu^*(\btheta,\rho)$ in probability.
By (A5), $\nabla_{\btheta}\sigma_{\cU_{\btheta,\rho}}(\bw_j)$ is continuous in $\rho$ near $\rho$,
so the envelope term and size-sensitivity term in \eqref{eq:ghat} are continuous in $\rho$.

\textbf{Step 4: Combine.}
Applying Slutsky's theorem to \eqref{eq:ghat}, combining Steps 1--3, yields
$\hat{\bg}_\varepsilon(\btheta)\to \nabla_{\btheta}J_{P,\varepsilon}(\btheta)$ in probability.
\end{proof}

\begin{theorem}[Coverage Preserved Under Tuned Training]\label{thm:coverage-preserved}
Under Assumption~\ref{asm:exchange} for the calibration sample and fresh $U_{\text{new}}$, let $\hat{\btheta}$ be computed using $\cD_{\text{train}}$ and $\cD_{\text{tune}}$ only (independent of $\cD_{\text{cal}}$). Compute $\hat{\rho}_\tau$ on $\cD_{\text{cal}}$ using~(\ref{eq:conformal-rho}) with shape $\hat{\btheta}$. Then
\begin{equation}
\Prob(U_{\text{new}} \in \cU_{\hat{\btheta}, \hat{\rho}_\tau}) \geq \tau.
\end{equation}
\end{theorem}

\begin{proof}
Condition on $\hat{\btheta}$. By construction, $\hat{\btheta}$ is independent of $\cD_{\text{cal}}$. Applying Theorem~\ref{thm:conformal} conditional on $\hat{\btheta}$ yields the result.
\end{proof}

\subsection{Ellipsoidal Specialization}

For power systems applications, uncertainty sets are parameterized as ellipsoids via Cholesky factors, which automatically ensure positive definiteness of the underlying covariance structure.

Let $\btheta = \bL \in \R^{d \times d}$ be a lower-triangular matrix with positive diagonal entries (Cholesky factor). The ellipsoidal gauge function is:
\begin{equation}
s_{\bL}(\bu) = \norm{\bL^{-1} \bu}_2.
\end{equation}
This yields the ellipsoid:
\begin{equation}
\cU_{\bL,\rho} = \{\bu : \bu^\top (\bL\bL^\top)^{-1} \bu \leq \rho^2\} = \{\bu : \norm{\bL^{-1}\bu}_2 \leq \rho\}.
\end{equation}

\begin{proposition}[Support Function for Ellipsoids]
For ellipsoid $\cU_{\bL,\rho}$:
\begin{equation}\label{eq:ellipsoid-support}
\sigma_{\cU_{\bL,\rho}}(\bw) = \rho \norm{\bL^\top \bw}_2.
\end{equation}
\end{proposition}

\begin{proof}
By change of variables $\bm{v} = \bL^{-1}\bu$:
\begin{align}
\sigma_{\cU_{\bL,\rho}}(\bw) &= \sup_{\norm{\bm{v}}_2 \leq \rho} \ip{\bw}{\bL\bm{v}} = \sup_{\norm{\bm{v}}_2 \leq \rho} \ip{\bL^\top\bw}{\bm{v}} = \rho \norm{\bL^\top\bw}_2.
\end{align}
\end{proof}

\begin{proposition}[Ellipsoidal Gradients]\label{prop:ellipsoid-grads}
Fix $\rho>0$. Consider the ellipsoidal gauge $s_{\bL}(\bu)=\|\bL^{-1}\bu\|_2$ and support function
$\sigma_{\cU_{\bL,\rho}}(\bw)=\rho\|\bL^\top\bw\|_2$.
All matrix gradients below are taken with respect to the Frobenius inner product
$\ipF{A}{B}:=\tr(A^\top B)$.

For $\bw\neq 0$ with $\bL^\top\bw\neq 0$,
\begin{equation}\label{eq:grad-sigma-L}
\nabla_{\bL}\,\sigma_{\cU_{\bL,\rho}}(\bw)
=\rho\,\frac{\bw(\bL^\top\bw)^\top}{\|\bL^\top\bw\|_2}
=\rho\,\frac{\bw\bw^\top\bL}{\|\bL^\top\bw\|_2}.
\end{equation}
For $\bw=0$, $\nabla_{\bL}\sigma_{\cU_{\bL,\rho}}(\bw)=0$.

For $\bu\neq 0$ with $\bL^{-1}\bu\neq 0$,
\begin{equation}\label{eq:grad-gauge-L}
\nabla_{\bL}\,s_{\bL}(\bu)
=-\,\frac{\bL^{-\top}(\bL^{-1}\bu)(\bL^{-1}\bu)^\top}{\|\bL^{-1}\bu\|_2}.
\end{equation}
At $\bu=0$, $s_{\bL}(\bu)$ is not differentiable, but $0$ is a valid choice of subgradient; in implementations we set $\nabla_{\bL}s_{\bL}(0)=0$.
\end{proposition}

\begin{proof}
Let $y(\bL):=\bL^\top\bw$. Then $\sigma_{\cU_{\bL,\rho}}(\bw)=\rho\|y(\bL)\|_2$.
For $\bL^\top\bw\neq 0$, the differential is
\[
\mathrm{d}\sigma
=\rho\,\frac{y^\top}{\|y\|_2}\,\mathrm{d}y
=\rho\,\frac{(\bL^\top\bw)^\top}{\|\bL^\top\bw\|_2}\,(\mathrm{d}\bL)^\top\bw
=\rho\,\ipF{\frac{\bw(\bL^\top\bw)^\top}{\|\bL^\top\bw\|_2}}{\mathrm{d}\bL}.
\]
This identifies the Frobenius gradient as \eqref{eq:grad-sigma-L}.
The $\bw=0$ case is immediate since $\sigma_{\cU_{\bL,\rho}}(0)=0$ for all $\bL$.

For the gauge, let $v(\bL):=\bL^{-1}\bu$. Then $s_{\bL}(\bu)=\|v(\bL)\|_2$.
For $v\neq 0$, $\mathrm{d}v = -\bL^{-1}(\mathrm{d}\bL)\bL^{-1}\bu = -\bL^{-1}(\mathrm{d}\bL)v$ and hence
\[
\mathrm{d}s
=\frac{v^\top}{\|v\|_2}\,\mathrm{d}v
=-\frac{v^\top}{\|v\|_2}\,\bL^{-1}(\mathrm{d}\bL)v
=-\ipF{\frac{\bL^{-\top}vv^\top}{\|v\|_2}}{\mathrm{d}\bL},
\]
which yields \eqref{eq:grad-gauge-L}.
\end{proof}

For numerical stability and identifiability, a trace normalization constraint is imposed:
\begin{equation}
\tr(\bL\bL^\top) = d.
\end{equation}
After each gradient step, $\bL$ is projected onto this constraint: $\bL \leftarrow \bL / \sqrt{\tr(\bL\bL^\top)/d}$.

\subsection{Context-Dependent Learning}

In practice, the correlation structure of forecast errors varies with system conditions. A neural network encoder $\bL_\phi: \cZ \to \R^{d \times d}$ maps context features $\bxi \in \cZ$ to Cholesky factors.

The contextual learning problem is:
\begin{equation}
\min_{\phi,\, \rho > 0} \E_{(\bxi,\bu) \sim P}\left[V(\bL_\phi(\bxi), \rho; \bxi)\right] \quad \text{s.t.} \quad \Prob(\bu \in \cU_{\bL_\phi(\bxi),\rho}) \geq \tau.
\end{equation}

Context features $\bxi$ include:
\begin{itemize}
\item \textbf{Weather:} temperature, wind speed, cloud cover, pressure gradient, humidity
\item \textbf{Temporal:} hour, day, month (cyclic encoding), weekday indicator, peak/ramp periods
\item \textbf{System state:} normalized load, renewable penetration, forecast horizon
\item \textbf{Economic:} zone-wise LMP-based reserve prices
\item \textbf{Lagged errors:} $\bu_{t-1}, \ldots, \bu_{t-k}$ for temporal correlation
\end{itemize}

The encoder architecture outputs a lower-triangular matrix:
\begin{enumerate}
\item Feature extraction: $\bm{h} = \text{MLP}_\phi(\bxi) \in \R^{d(d+1)/2}$
\item Cholesky construction: Fill lower triangle with $\bm{h}$, apply $\exp(\cdot)$ to diagonal for positivity
\item Normalization: $\bL \leftarrow \bL / \sqrt{\tr(\bL\bL^\top)/d}$
\end{enumerate}

\subsection{Two Structural Regimes}

The methodology admits two structural regimes depending on the coupling structure of the robust problem.

\textbf{Regime 1: Zonal Reserve (Closed-Form Reserve Requirements).}
When reserve constraints are separable across zones and transmission constraints are non-binding,
the robust reserve requirements reduce to closed-form expressions.
For zone $z$ with exposure vector $\bm{e}_z$:
\begin{equation}
R_z^{\min}=\sigma_{\cU_{\bL,\rho}}(\bm{e}_z)=\rho\|\bL^\top \bm{e}_z\|_2
=\rho\sqrt{\sum_{z'=1}^d L_{z'z}^2}.
\end{equation}
In this regime:
\begin{itemize}
\item The reserve requirements are explicit functions of $(\bL,\rho)$, hence their gradients are closed-form via Proposition~\ref{prop:ellipsoid-grads}.
\item Gradient computation does not require differentiating through (or even solving) a network-constrained SCED.
\item If a (separable) reserve procurement subproblem is solved, its dual multipliers can be recovered from its KKT conditions; however, they are model-dependent and are not assumed to equal cost coefficients in general.
\end{itemize}

\textbf{Regime 2: Network-Constrained (Envelope Gradient).} When transmission limits are binding, the SCED problem must be solved to obtain optimal dispatch and dual multipliers. The gradient is computed via the envelope theorem:
\begin{equation}
\nabla_{\bL} V = \sum_j \mu_j^* \nabla_{\bL} \sigma_{\cU_{\bL,\rho}}(\bw_j)
\end{equation}
where $\mu_j^*$ are extracted from the SCED solution.

%===============================================================================
% IV. IEEE 118-Bus System Formulation
%===============================================================================
\section{IEEE 118-Bus System Formulation}
\label{sec:case-study}

This section instantiates the general framework for the IEEE~118-bus test system, a standard benchmark for power system studies.

\subsection{Network Description}

\begin{table}[!t]
\caption{IEEE 118-Bus System Parameters}
\label{tab:system-params}
\centering
\begin{tabular}{ll}
\toprule
\textbf{Parameter} & \textbf{Value} \\
\midrule
Buses & 118 \\
Generators & 54 \\
Transmission Lines & 186 \\
Load Zones & 10 (aggregated) \\
Uncertainty Dimension & $d = 10$ \\
Total Load (base case) & 4,242 MW \\
Generation Capacity & $\sim$5,500 MW \\
\bottomrule
\end{tabular}
\end{table}

Table~\ref{tab:system-params} summarizes the system parameters. Buses are aggregated into $d = 10$ load zones for uncertainty modeling, with each zone's net load forecast error represented as a component of the uncertainty vector $\bu \in \R^{10}$.

\begin{table}[!t]
\caption{Zonal Aggregation}
\label{tab:zones}
\centering
\begin{tabular}{cccc}
\toprule
\textbf{Zone} & \textbf{Buses} & \textbf{Load (MW)} & \textbf{Gen. Cap. (MW)} \\
\midrule
1 & 1--12 & 423 & 550 \\
2 & 13--24 & 412 & 520 \\
3 & 25--36 & 445 & 580 \\
4 & 37--48 & 398 & 490 \\
5 & 49--60 & 467 & 610 \\
6 & 61--72 & 389 & 480 \\
7 & 73--84 & 456 & 590 \\
8 & 85--96 & 401 & 510 \\
9 & 97--108 & 478 & 620 \\
10 & 109--118 & 373 & 550 \\
\midrule
\textbf{Total} & \textbf{118} & \textbf{4,242} & \textbf{5,500} \\
\bottomrule
\end{tabular}
\end{table}

\subsection{Security-Constrained Economic Dispatch}

The decision variables are:
\begin{itemize}
\item $g_i \in \R$: Generation dispatch for generator $i$ (MW)
\item $r_i \in \R_+$: Reserve allocation for generator $i$ (MW)
\item $\theta_b \in \R$: Voltage angle at bus $b$ (rad)
\end{itemize}

The SCED problem minimizes total generation and reserve costs:
\begin{equation}\label{eq:sced-obj}
\min_{g, r, \theta} \sum_{i \in \cG} \left(c_i^g g_i + c_i^r r_i\right)
\end{equation}
where $c_i^g$ is the generation cost (\$/MWh) and $c_i^r$ is the reserve cost (\$/MWh).

\textbf{Power Balance at Each Bus:}
\begin{equation}\label{eq:power-balance}
\sum_{i \in \cG_b} g_i - D_b = \sum_{(b,b') \in \cL} B_{bb'}(\theta_b - \theta_{b'}), \quad \forall b \in \cB
\end{equation}
where $B_{bb'} = 1/X_{bb'}$ is the line susceptance and $D_b$ is the load at bus $b$.

\textbf{DC Power Flow Line Limits:}
\begin{equation}\label{eq:line-limits}
-F_\ell^{\max} \leq B_\ell(\theta_{b(\ell)} - \theta_{b'(\ell)}) \leq F_\ell^{\max}, \quad \forall \ell \in \cL
\end{equation}
where $F_\ell^{\max}$ is the thermal limit of line $\ell$.

\textbf{Generator Capacity:}
\begin{equation}\label{eq:gen-capacity}
\underline{g}_i \leq g_i \leq \bar{g}_i - r_i, \quad r_i \geq 0, \quad \forall i \in \cG
\end{equation}

\textbf{Zonal Reserve Requirements:}
\begin{equation}\label{eq:zonal-reserve}
\sum_{i \in \cG_z} r_i \geq R_z^{\min}, \quad \forall z \in \cZ
\end{equation}
where the minimum reserve requirement for zone $z$ is:
\begin{equation}\label{eq:reserve-requirement}
R_z^{\min} = \sigma_{\cU_{\bL,\rho}}(\bm{e}_z) = \rho \norm{\bL^\top \bm{e}_z}_2 = \rho \sqrt{\sum_{z'=1}^{d} L_{z'z}^2}.
\end{equation}

\subsection{PTDF-Based Network-Constrained Formulation}

We use a PTDF representation of DC power flow (equivalent to the angle formulation with a fixed slack bus).
Let $\cB$ denote the buses, $\cL$ the lines, and $\cG$ the generators.

\textbf{Bus injections from generator variables.}
Define the net injection at bus $b$ under dispatch $g$ as
\[
p_b(g) := \sum_{i\in\cG_b} g_i - D_b,
\qquad b\in\cB,
\]
and let $p(g)\in\R^{|\cB|}$ collect these injections.

\textbf{PTDF mapping.}
Let $\Pi\in\R^{|\cL|\times|\cB|}$ be the PTDF matrix mapping bus injections to line flows:
\[
f(g) := \Pi\,p(g)\in\R^{|\cL|},
\qquad
f_\ell(g) = \sum_{b\in\cB}\Pi_{\ell b}\,p_b(g).
\]

\textbf{Zonal uncertainty and line-flow exposure.}
Zonal net-load forecast error is modeled by $\bu\in\R^d$. Let $\Pi^{\text{zone}}\in\R^{|\cL|\times d}$ denote the linear map from zonal errors to line-flow deviations.
Thus line $\ell$ experiences an uncertain deviation $\ip{\Pi_\ell^{\text{zone}}}{\bu}$, where $\Pi_\ell^{\text{zone}}\in\R^d$ is row $\ell$.

\begin{lemma}[Robust Counterpart of Absolute-Value Constraints]\label{lem:robust-abs}
Let $\cU\subseteq\R^d$ be nonempty and let $\bw\in\R^d$. For any scalar $f\in\R$,
\[
\sup_{\bu\in\cU} \big| f + \ip{\bw}{\bu} \big|
\le F
\]
is equivalent to the pair of inequalities
\begin{equation}\label{eq:robust-abs-pair}
f + \sigma_{\cU}(\bw)\le F,
\qquad
-f + \sigma_{\cU}(-\bw)\le F.
\end{equation}
If $\cU$ is centrally symmetric (i.e., $\cU=-\cU$), then $\sigma_{\cU}(-\bw)=\sigma_{\cU}(\bw)$ and
\eqref{eq:robust-abs-pair} is equivalent to
\[
|f| + \sigma_{\cU}(\bw)\le F.
\]
\end{lemma}

\begin{proof}
The robust inequality is equivalent to the conjunction of
$f+\ip{\bw}{\bu}\le F$ for all $\bu\in\cU$ and $-(f+\ip{\bw}{\bu})\le F$ for all $\bu\in\cU$.
Taking suprema over $\bu\in\cU$ yields \eqref{eq:robust-abs-pair}.
If $\cU=-\cU$, then $\sigma_{\cU}(-\bw)=\sup_{\bu\in\cU}\ip{-\bw}{\bu}=\sup_{\bu\in\cU}\ip{\bw}{-\bu}=\sigma_{\cU}(\bw)$,
and the pair reduces to $f+\sigma_{\cU}(\bw)\le F$ and $-f+\sigma_{\cU}(\bw)\le F$, i.e., $|f|+\sigma_{\cU}(\bw)\le F$.
\end{proof}

\textbf{Robust line limits.}
We enforce
\[
\big| f_\ell(g) + \ip{\Pi_\ell^{\text{zone}}}{\bu} \big|\le F_\ell^{\max}\qquad \forall \bu\in\cU_{\btheta,\rho}.
\]
By Lemma~\ref{lem:robust-abs}, the exact robust counterpart for general $\cU_{\btheta,\rho}$ is the pair
\begin{equation}\label{eq:robust-line}
\begin{aligned}
f_\ell(g) + \sigma_{\cU_{\btheta,\rho}}(\Pi_\ell^{\text{zone}}) &\le F_\ell^{\max},\\
-f_\ell(g) + \sigma_{\cU_{\btheta,\rho}}(-\Pi_\ell^{\text{zone}}) &\le F_\ell^{\max},
\end{aligned}
\qquad \forall \ell\in\cL.
\end{equation}
In the ellipsoidal specialization $\cU_{\bL,\rho}$ (Section~\ref{sec:methodology}-E), the set is centrally symmetric, so
$\sigma_{\cU_{\bL,\rho}}(-\bw)=\sigma_{\cU_{\bL,\rho}}(\bw)$ and \eqref{eq:robust-line} reduces to
\[
|f_\ell(g)| + \sigma_{\cU_{\bL,\rho}}(\Pi_\ell^{\text{zone}})\le F_\ell^{\max}
= |f_\ell(g)| + \rho \norm{\bL^\top \Pi_\ell^{\text{zone}}}_2 \le F_\ell^{\max}.
\]

\textbf{Network-constrained robust SCED.}
Using the exact robust counterpart \eqref{eq:robust-line}, the network-constrained problem is
\begin{equation}\label{eq:network-sced}
\begin{aligned}
\min_{g,r}\quad & \sum_{i \in \cG}\big(c_i^g g_i + c_i^r r_i\big) \\
\text{s.t.}\quad
& \sum_{i\in\cG} g_i = \sum_{b\in\cB} D_b, \\
& f_\ell(g) + \sigma_{\cU_{\btheta,\rho}}(\Pi_\ell^{\text{zone}}) \le F_\ell^{\max}, \quad \forall \ell\in\cL,\\
& -f_\ell(g) + \sigma_{\cU_{\btheta,\rho}}(-\Pi_\ell^{\text{zone}}) \le F_\ell^{\max}, \quad \forall \ell\in\cL,\\
& \underline g_i \le g_i \le \bar g_i - r_i,\quad r_i\ge 0,\quad \forall i\in\cG,\\
& \sum_{i \in \cG_z} r_i \ge \sigma_{\cU_{\btheta,\rho}}(\bm{e}_z), \quad \forall z\in\cZ,
\end{aligned}
\end{equation}
where $f_\ell(g) = \sum_{b\in\cB}\Pi_{\ell b}\,p_b(g)$ is the line flow under dispatch $g$.

\subsection{Uncertainty Model}

The uncertainty vector $\bu \in \R^d$ represents zonal net load forecast errors:
\begin{equation}
u_z = (\text{Actual Load}_z - \text{Forecast Load}_z) - (\text{Actual RE}_z - \text{Forecast RE}_z)
\end{equation}
where RE denotes renewable energy generation.

Empirical studies indicate that renewable forecast errors exhibit:
\begin{itemize}
\item \textbf{Spatial correlation:} Adjacent zones have correlation $\rho_{zz'} \approx 0.3$--$0.7$ due to common weather patterns
\item \textbf{Temporal correlation:} Lag-1 autocorrelation $\approx 0.5$--$0.8$
\item \textbf{Weather dependence:} Correlation increases during weather events
\end{itemize}

The learned Cholesky factor $\bL$ captures these correlation patterns adaptively based on context~$\bxi$.

%===============================================================================
% V. Algorithm
%===============================================================================
\section{Algorithm}
\label{sec:algorithm}

Algorithm~\ref{alg:main} presents the complete training procedure with tuning-based profiled gradient approximation and final conformal calibration.

\begin{algorithm}[!t]
\caption{Tuning-Based Profiled Gradient Training with Conformal Calibration}
\label{alg:main}
\begin{algorithmic}[1]
\REQUIRE Training data $\cD_{\text{train}}$, tuning data $\cD_{\text{tune}}$, calibration data $\cD_{\text{cal}}$, coverage level $\tau$, bandwidth $\varepsilon$, iterations $K$, step sizes $\{\eta_k\}$
\STATE Initialize shape parameter $\btheta_0 \in \Theta$

\FOR{$k = 0, 1, \ldots, K-1$}
    \STATE \textbf{// Tuning: estimate radius and quantile sensitivity}
    \STATE Compute gauge scores $S_i(\btheta_k) = s_{\btheta_k}(U_i)$ for $U_i \in \cD_{\text{tune}}$
    \STATE Compute smoothed $\tau$-quantile: $\hat{\rho}_\varepsilon(\btheta_k)$
    \STATE Compute weights: $\omega_i(\btheta_k) = \varphi((\hat{\rho}_\varepsilon(\btheta_k) - S_i(\btheta_k))/\varepsilon)$
    \STATE Estimate quantile gradient: $\widehat{\nabla_{\btheta}\rho_\varepsilon}(\btheta_k)$ via~(\ref{eq:rhohat-grad})

    \STATE \textbf{// Robust solve: obtain dual multipliers}
    \STATE Solve robust problem~(\ref{eq:robust}) at $(\btheta_k, \hat{\rho}_\varepsilon(\btheta_k))$
    \STATE Extract optimal dual multipliers $\bmu^*(\btheta_k, \hat{\rho}_\varepsilon(\btheta_k))$

    \STATE \textbf{// Compute approximate profiled gradient}
    \STATE Envelope shape term: $\bg_{\text{shape}} = \sum_j \mu_j^* \nabla_{\btheta} \sigma_j(\btheta_k, \hat{\rho}_\varepsilon(\btheta_k))$
    \STATE Envelope size sensitivity: $g_{\text{size}} = \sum_j \mu_j^* \sigma_{\cU_{\btheta_k,1}}(\bw_j)$
    \STATE Full gradient: $\hat{\bg}_\varepsilon(\btheta_k) = \bg_{\text{shape}} + g_{\text{size}} \cdot \widehat{\nabla_{\btheta}\rho_\varepsilon}(\btheta_k)$

    \STATE \textbf{// Update}
    \STATE $\btheta_{k+1} \leftarrow \Pi_\Theta(\btheta_k - \eta_k \cdot \hat{\bg}_\varepsilon(\btheta_k))$
\ENDFOR

\STATE \textbf{// Final conformal calibration (independent of training)}
\STATE Set final shape: $\hat{\btheta} = \btheta_K$
\STATE Compute calibration scores: $S_i = s_{\hat{\btheta}}(U_i)$ for $U_i \in \cD_{\text{cal}}$
\STATE Sort: $S_{(1)} \leq S_{(2)} \leq \cdots \leq S_{(n_{\text{cal}})}$
\STATE Set $k = \lceil(n_{\text{cal}}+1)\tau\rceil$
\STATE Calibrated radius: $\hat{\rho}_\tau = S_{(k)}$

\ENSURE Learned shape $\hat{\btheta}$, calibrated radius $\hat{\rho}_\tau$ with coverage guarantee
\end{algorithmic}
\end{algorithm}

\begin{remark}[Computational Complexity]
In the solver-free zonal regime (Regime~1), the algorithm requires no optimization solver calls during training---all computations are closed-form. In the network-constrained regime (Regime~2), one SCED solve per iteration is required to obtain dual multipliers, but no differentiation through the solver is needed.
\end{remark}

%===============================================================================
% VI. Discussion
%===============================================================================
\section{Discussion}
\label{sec:discussion}

\subsection{Practical Implications for Market Design}

The proposed framework has several implications for electricity market design:

\begin{enumerate}
\item \textbf{Reserve Procurement Efficiency:} Learning correlation structure enables more accurate reserve requirements, reducing over-procurement while maintaining reliability.

\item \textbf{Price Signal Accuracy:} Better uncertainty modeling leads to more accurate locational marginal prices (LMPs), as reserve costs are more accurately allocated to zones.

\item \textbf{Renewable Integration:} Context-dependent uncertainty sets adapt to varying correlation patterns as renewable penetration increases, supporting efficient integration.
\end{enumerate}

\subsection{Theoretical Insights on Correlation}

The effect of correlation structure on reserve requirements can be analyzed theoretically.

\textbf{Two-Zone Example.} Consider a two-zone system with covariance
\begin{equation}
\Sigma = \begin{pmatrix} \sigma_1^2 & \rho_{12}\sigma_1\sigma_2 \\ \rho_{12}\sigma_1\sigma_2 & \sigma_2^2 \end{pmatrix}
\end{equation}
where $\rho_{12} \in [-1,1]$ is the correlation coefficient.

For zonal reserve requirements (sum of independent constraints):
\begin{equation}
R_{\text{zonal}} = \sqrt{\bm{e}_1^\top \Sigma \bm{e}_1} + \sqrt{\bm{e}_2^\top \Sigma \bm{e}_2} = \sigma_1 + \sigma_2.
\end{equation}

For joint reserve (system-wide constraint on $\bu_1 + \bu_2$):
\begin{equation}
R_{\text{joint}} = \sqrt{(\bm{e}_1 + \bm{e}_2)^\top \Sigma (\bm{e}_1 + \bm{e}_2)} = \sqrt{\sigma_1^2 + \sigma_2^2 + 2\rho_{12}\sigma_1\sigma_2}.
\end{equation}

\textbf{Key Insight:}
\begin{itemize}
\item If $\rho_{12} > 0$: $R_{\text{joint}} > \sqrt{\sigma_1^2 + \sigma_2^2}$ (risk amplification)
\item If $\rho_{12} < 0$: $R_{\text{joint}} < \sqrt{\sigma_1^2 + \sigma_2^2}$ (risk hedging)
\item If $\rho_{12} = 1$: $R_{\text{joint}} = \sigma_1 + \sigma_2$ (maximum correlation)
\item If $\rho_{12} = -1$ and $\sigma_1 = \sigma_2$: $R_{\text{joint}} = 0$ (perfect hedge)
\end{itemize}

\textbf{General $d$-Zone System with Equicorrelation.} For equal variance $\sigma^2$ and uniform correlation $\rho$ across all zones:
\begin{equation}
\Sigma = \sigma^2\left[(1-\rho)I_d + \rho \bm{1}\bm{1}^\top\right].
\end{equation}

The system-level reserve for a joint constraint is:
\begin{equation}
R_{\text{sys}} = \sigma\sqrt{d(1-\rho) + d^2\rho} = \sigma\sqrt{d}\sqrt{1 + (d-1)\rho}.
\end{equation}

For $d = 10$ zones:
\begin{itemize}
\item $\rho = 0$: $R_{\text{sys}} = \sigma\sqrt{10} \approx 3.16\sigma$
\item $\rho = 0.5$: $R_{\text{sys}} = \sigma\sqrt{55} \approx 7.42\sigma$
\item $\rho = 0.9$: $R_{\text{sys}} = \sigma\sqrt{91} \approx 9.54\sigma$
\end{itemize}

This demonstrates that positive correlation substantially increases system reserve requirements, while learning the true correlation structure (rather than assuming independence) enables appropriate reserve sizing.

\subsection{Computational Considerations}

\begin{table}[!t]
\caption{Comparison of Two Structural Regimes}
\label{tab:regimes}
\centering
\begin{tabular}{lcc}
\toprule
& \textbf{Regime 1} & \textbf{Regime 2} \\
& \textbf{(Zonal)} & \textbf{(Network)} \\
\midrule
Reserve constraints & Separable & Coupled \\
Transmission limits & Non-binding & Binding \\
Gradient computation & Closed-form & Envelope \\
Solver required & No & Yes (per iteration) \\
Dual multipliers & From KKT (if solved) & From SCED \\
Computational cost & Low & Moderate \\
\bottomrule
\end{tabular}
\end{table}

Table~\ref{tab:regimes} compares the two structural regimes. The solver-free regime enables efficient training when transmission constraints are non-binding, while the envelope gradient approach handles general network-constrained problems without differentiating through the solver.

\subsection{Limitations and Extensions}

\begin{enumerate}
\item \textbf{Data Requirements:} The framework requires sufficient historical data to learn correlation patterns. Performance may degrade with limited data or in rapidly changing systems.

\item \textbf{Extreme Events:} Learned models may not generalize well to extreme events outside the training distribution. Robust extensions incorporating worst-case bounds merit investigation.

\item \textbf{Multi-Period Formulations:} Extension to multi-period unit commitment with ramping constraints and inter-temporal coupling is a natural direction.

\item \textbf{AC Power Flow:} The current formulation uses DC power flow approximation. Extension to AC optimal power flow requires additional technical development.
\end{enumerate}

%===============================================================================
% VII. Conclusion
%===============================================================================
\section{Conclusion}
\label{sec:conclusion}

This paper developed a differentiable robust optimization framework for learning correlated uncertainty sets in operating reserve markets. The key innovations are:

\begin{enumerate}
\item \textbf{Gauge-based uncertainty sets} parameterized by shape $\btheta$ and size $\rho$, enabling flexible geometry learning while maintaining convexity.

\item \textbf{Profiled gradient derivation} that accounts for both direct shape effects and quantile-sensitivity corrections, providing the true gradient of the coverage-constrained objective.

\item \textbf{Tuning-based gradient approximation} using smoothed CDF estimation to approximate the profiled gradient when the distribution is unknown.

\item \textbf{Three-way data split} (train/tune/calibrate) that preserves distribution-free conformal coverage guarantees while enabling gradient-based shape optimization.

\item \textbf{Envelope-based gradients} that avoid differentiation through optimization solvers, using only dual multipliers from the robust problem.

\item \textbf{Two structural regimes} with closed-form gradients (zonal reserve) and envelope gradients (network-constrained), providing computational efficiency appropriate to problem structure.
\end{enumerate}

The framework was instantiated for security-constrained economic dispatch on the IEEE~118-bus system, demonstrating how learned correlation structure affects reserve procurement. Theoretical analysis shows that positive correlation increases system reserve requirements while negative correlation enables hedging, highlighting the importance of accurate correlation modeling for efficient market operations.

Future work includes extension to multi-period formulations, integration with real-time balancing markets, and development of robust extensions for extreme event scenarios.

%===============================================================================
% Appendix
%===============================================================================
\appendices

\section{Proof of Envelope Theorem}
\label{app:envelope-proof}

This appendix provides the complete proof of Theorem~\ref{thm:envelope-theta}.

\begin{proof}[Proof of Theorem~\ref{thm:envelope-theta}]
Define the Lagrangian for problem~(\ref{eq:robust}):
\begin{equation}
\cL(\bx, \bmu; \btheta, \rho) = f(\bx) + I_{\cX}(\bx) + \sum_{j=1}^m \mu_j(a_j(\bx) + \sigma_j(\btheta,\rho) - b_j)
\end{equation}
where $I_{\cX}(\bx) = 0$ if $\bx \in \cX$ and $+\infty$ otherwise.

The dual function is:
\begin{equation}
q(\bmu; \btheta, \rho) = \inf_{\bx \in \R^{n_x}} \cL(\bx, \bmu; \btheta, \rho).
\end{equation}

Since $\sigma_j(\btheta,\rho)$ does not depend on $\bx$, we can separate:
\begin{align}
q(\bmu; \btheta, \rho) &= \inf_{\bx}\left\{f(\bx) + I_{\cX}(\bx) + \sum_{j=1}^m \mu_j a_j(\bx)\right\} \nonumber \\
&\quad + \sum_{j=1}^m \mu_j(\sigma_j(\btheta,\rho) - b_j) \nonumber \\
&= \tilde{q}(\bmu) + \sum_{j=1}^m \mu_j(\sigma_j(\btheta,\rho) - b_j)
\end{align}
where $\tilde{q}(\bmu)$ is independent of $(\btheta, \rho)$.

For fixed $(\rho, \bmu)$, the dual function is differentiable in $\btheta$ with:
\begin{equation}
\nabla_{\btheta} q(\bmu; \btheta, \rho) = \sum_{j=1}^m \mu_j \nabla_{\btheta} \sigma_j(\btheta, \rho).
\end{equation}

By Assumption~\ref{asm:slater}, strong duality holds:
\begin{equation}
V(\btheta, \rho) = \max_{\bmu \geq 0} q(\bmu; \btheta, \rho)
\end{equation}
and the maximum is attained at some $\bmu^*(\btheta, \rho)$.

\textbf{Clarke subgradient property.} Fix $\btheta$ and let $\bmu^* \in \cM^*(\btheta, \rho)$. For any $\btheta'$:
\begin{align}
V(\btheta', \rho) &= \max_{\bmu \geq 0} q(\bmu; \btheta', \rho) \geq q(\bmu^*; \btheta', \rho).
\end{align}

By differentiability of $q(\bmu^*; \cdot, \rho)$ at $\btheta$:
\begin{align}
q(\bmu^*; \btheta', \rho) &= q(\bmu^*; \btheta, \rho) + \ip{\nabla_{\btheta} q(\bmu^*; \btheta, \rho)}{\btheta' - \btheta} \nonumber \\
&\quad + o(\norm{\btheta' - \btheta}).
\end{align}

Since $q(\bmu^*; \btheta, \rho) = V(\btheta, \rho)$ by optimality of $\bmu^*$:
\begin{align}
V(\btheta', \rho) - V(\btheta, \rho) &\geq \ip{\nabla_{\btheta} q(\bmu^*; \btheta, \rho)}{\btheta' - \btheta} \nonumber \\
&\quad + o(\norm{\btheta' - \btheta}).
\end{align}

This establishes that $\nabla_{\btheta} q(\bmu^*; \btheta, \rho) = \sum_j \mu_j^* \nabla_{\btheta} \sigma_j(\btheta, \rho) = \bg_{\btheta}(\btheta, \rho;\bmu^*)$ is a Clarke subgradient of $V(\cdot, \rho)$ at $\btheta$, i.e., $\bg_{\btheta}(\btheta,\rho;\bmu^*) \in \clsubd_{\btheta} V(\btheta,\rho)$.

\textbf{Differentiability.} If the dual maximizer $\bmu^*(\btheta, \rho)$ is unique (i.e., $\cM^*(\btheta,\rho) = \{\bmu^*\}$), then by standard Danskin/envelope theorem arguments, the pointwise maximum $V(\btheta, \rho) = \max_{\bmu \geq 0} q(\bmu; \btheta, \rho)$ is differentiable at $\btheta$, and the Clarke subdifferential reduces to a singleton equal to the gradient.
\end{proof}

\section{Conformal Coverage Proof}
\label{app:conformal-proof}

This appendix provides the complete proof of Theorem~\ref{thm:conformal}.

\begin{proof}[Proof of Theorem~\ref{thm:conformal}]
Consider the $(n_{\text{cal}} + 1)$ scores:
\begin{equation}
S_1, S_2, \ldots, S_{n_{\text{cal}}}, S_{\text{new}}
\end{equation}
where $S_i = s_{\hat{\btheta}}(U_i)$ for $i = 1, \ldots, n_{\text{cal}}$ and $S_{\text{new}} = s_{\hat{\btheta}}(U_{\text{new}})$.

\textbf{Exchangeability of scores.} By Assumption~\ref{asm:exchange}, $(U_1, \ldots, U_{n_{\text{cal}}}, U_{\text{new}})$ are exchangeable. Since $\hat{\btheta}$ is independent of all these random variables, the transformation $U \mapsto s_{\hat{\btheta}}(U)$ preserves exchangeability. Thus $(S_1, \ldots, S_{n_{\text{cal}}}, S_{\text{new}})$ are exchangeable.

\textbf{Rank uniformity.} For exchangeable random variables, each variable is equally likely to have any rank among the group. Specifically, the rank of $S_{\text{new}}$ among $\{S_1, \ldots, S_{n_{\text{cal}}}, S_{\text{new}}\}$ is uniformly distributed on $\{1, 2, \ldots, n_{\text{cal}} + 1\}$.

(In the case of ties, we can use a consistent tie-breaking rule or note that continuous score distributions have ties with probability zero.)

\textbf{Coverage event.} Let $S_{(1)} \leq S_{(2)} \leq \cdots \leq S_{(n_{\text{cal}})}$ be the order statistics of $S_1, \ldots, S_{n_{\text{cal}}}$. The event $S_{\text{new}} \leq S_{(k)}$ occurs if and only if the rank of $S_{\text{new}}$ among all $(n_{\text{cal}} + 1)$ scores is at most $k$.

\textbf{Coverage probability.}
\begin{align}
\Prob(S_{\text{new}} \leq S_{(k)}) &= \Prob(\text{rank}(S_{\text{new}}) \leq k) \nonumber \\
&= \frac{k}{n_{\text{cal}} + 1} \nonumber \\
&\geq \frac{\lceil (n_{\text{cal}} + 1)\tau \rceil}{n_{\text{cal}} + 1} \nonumber \\
&\geq \tau.
\end{align}

The last inequality holds because $\lceil x \rceil \geq x$ for all $x \in \R$.

Since $S_{\text{new}} \leq \hat{\rho}_\tau = S_{(k)}$ is equivalent to $U_{\text{new}} \in \cU_{\hat{\btheta}, \hat{\rho}_\tau}$, the coverage guarantee follows.
\end{proof}

%===============================================================================
% References
%===============================================================================
\begin{thebibliography}{99}

\bibitem{Hodge2011}
B.-M. Hodge and M. Milligan, ``Wind power forecasting error distributions over multiple timescales,'' in \emph{Proc. IEEE Power Energy Soc. General Meeting}, Detroit, MI, USA, 2011, pp.~1--8.

\bibitem{Pinson2013}
P. Pinson, ``Wind energy: Forecasting challenges for its operational management,'' \emph{Statist. Sci.}, vol.~28, no.~4, pp.~564--585, 2013.

\bibitem{BenTal2009}
A. Ben-Tal, L. El Ghaoui, and A. Nemirovski, \emph{Robust Optimization}. Princeton, NJ, USA: Princeton Univ. Press, 2009.

\bibitem{Bertsimas2011}
D. Bertsimas, D. B. Brown, and C. Caramanis, ``Theory and applications of robust optimization,'' \emph{SIAM Rev.}, vol.~53, no.~3, pp.~464--501, 2011.

\bibitem{Bertsimas2013}
D. Bertsimas, E. Litvinov, X. A. Sun, J. Zhao, and T. Zheng, ``Adaptive robust optimization for the security constrained unit commitment problem,'' \emph{IEEE Trans. Power Syst.}, vol.~28, no.~1, pp.~52--63, Feb. 2013.

\bibitem{Jiang2012}
R. Jiang, J. Wang, and Y. Guan, ``Robust unit commitment with wind power and pumped storage hydro,'' \emph{IEEE Trans. Power Syst.}, vol.~27, no.~2, pp.~800--810, May 2012.

\bibitem{Zeng2013}
B. Zeng and L. Zhao, ``Solving two-stage robust optimization problems using a column-and-constraint generation method,'' \emph{Oper. Res. Lett.}, vol.~41, no.~5, pp.~457--461, 2013.

\bibitem{Bertsimas2018}
D. Bertsimas, V. Gupta, and N. Kallus, ``Data-driven robust optimization,'' \emph{Math. Program.}, vol.~167, no.~2, pp.~235--292, 2018.

\bibitem{Esfahani2018}
P. M. Esfahani and D. Kuhn, ``Data-driven distributionally robust optimization using the Wasserstein metric: Performance guarantees and tractable reformulations,'' \emph{Math. Program.}, vol.~171, no.~1-2, pp.~115--166, 2018.

\bibitem{Amos2017}
B. Amos and J. Z. Kolter, ``OptNet: Differentiable optimization as a layer in neural networks,'' in \emph{Proc. Int. Conf. Mach. Learn.}, Sydney, Australia, 2017, pp.~136--145.

\bibitem{Agrawal2019}
A. Agrawal, B. Amos, S. Barratt, S. Boyd, S. Diamond, and J. Z. Kolter, ``Differentiable convex optimization layers,'' in \emph{Proc. Adv. Neural Inf. Process. Syst.}, Vancouver, Canada, 2019, pp.~9562--9574.

\bibitem{Vovk2005}
V. Vovk, A. Gammerman, and G. Shafer, \emph{Algorithmic Learning in a Random World}. New York, NY, USA: Springer, 2005.

\bibitem{Romano2019}
Y. Romano, E. Patterson, and E. Cand\`{e}s, ``Conformalized quantile regression,'' in \emph{Proc. Adv. Neural Inf. Process. Syst.}, Vancouver, Canada, 2019, pp.~3543--3553.

\bibitem{Tibshirani2019}
R. J. Tibshirani, R. Foygel Barber, E. Cand\`{e}s, and A. Ramdas, ``Conformal prediction under covariate shift,'' in \emph{Proc. Adv. Neural Inf. Process. Syst.}, Vancouver, Canada, 2019, pp.~2530--2540.

\bibitem{IEEE118}
IEEE 118-bus test system data. [Online]. Available: \url{https://icseg.iti.illinois.edu/ieee-118-bus-system/}

\end{thebibliography}

\end{document}
